% 重要符号表（照录原书 pp.i–iv；转录自 300dpi 页 p-007–p-010；按原书四页分四个 tabular）
% ---------- 原书第 i 页 ----------
{\small
\begin{tabular}{@{}p{0.46\textwidth}@{}p{0.50\textwidth}@{}}
$A$ & 矢量势 \\
$\langle A\rangle$ & 统计平均 \\
$a_{\rm H}$ & 玻尔半径 \\
$a_{\rm ex}$ & 激子半径 \\
$\ve{a}_i\ (i=1,2,3)$ & 正点阵基矢 \\
$\ve{b}_i\ (i=1,2,3)$ & 倒点阵基矢 \\
BZ & 第一布里渊区（简约区） \\
$a^+,\,(b^+)$ & 玻色子产生算符 \\
$a,\,(b)$ & 玻色子消灭算符 \\
$C^+(C)$ & 费密子产生（消灭）算符 \\
$c_{\rm L}\,(c_{\rm T})$ & 纵（横）波声速 \\
$c_{\rm es}$ & 超导电子比热容 \\
$c_{\rm m}$ & 磁振子比热容 \\
$\ve{D}$ & 电位移矢量 \\
$D$ & 扩散系数 \\
$D_{\alpha\beta}(\ve{k}),\ D_{\alpha\beta}(\ve{k}_{s,s'})$ & 动力矩阵 \\
$e$ & 电子电荷 \\
$\ve{e}_{k\sigma}$ & 格波的极化矢量 \\
$E_{\rm F}$ & 费密能 \\
$E_{\rm ex}$ & 交换能 \\
$E_{\rm corr}$ & 相关能 \\
$E_{\rm k}$ & 动能 \\
$E_n(-{\rm i}\nabla),\ \mathscr{E}(-{\rm i}\nabla)$ & 有效哈密顿量 \\
$F$ & 自由能 \\
$f$ & 费密分布 \\
$G$ & 吉布斯函数 \\
$g$ & 朗德因子 \\
$g(E)$ & 一种自旋取向的电子态密度 \\
$G_{\rm r}(t-t')\,[\,G_{\rm a}(t-t')\,]$ & 双时推迟（超前）格林函数 \\
$G_{\rm r}(\omega)\equiv\langle\langle A|B\rangle\rangle_{\omega+{\rm i}\eta},\ \eta=+0$ & 双时推迟格林函数的傅里叶变换 \\
\end{tabular}
}
\par\bigbreak
% ---------- 原书第 ii 页 ----------
{\small
\begin{tabular}{@{}p{0.46\textwidth}@{}p{0.50\textwidth}@{}}
$G_{\rm a}(\omega)\equiv\langle\langle A|B\rangle\rangle_{\omega-{\rm i}\eta},\ \eta=+0$ & 双时超前格林函数的傅里叶变换 \\
$H,\ \mathscr{H}$ & 哈密顿量 \\
$\overline{H}$ & BCS 约化哈密顿量 \\
$h$ & 普朗克常量 \\
$\hbar$ & $\hbar=h/2\pi$ \\
$I\,(I_{NN},I_{SS})$ & 隧道电流 \\
$\ve{j}\,(\ve{r})$ & 电流密度 \\
$\ve{K}\equiv\ve{K}_n=\sum_{i=1}^{s}n_i\ve{b}_i$ & 倒格矢 \\
$\ve{k}$ & 波矢量 \\
$k_{\rm B}$ & 玻耳兹曼常量 \\
$k_{\rm F}$ & 费密波数 \\
$\mathscr{L}$ & 拉氏量密度 \\
$M$ & 原子（分子）质量 \\
$\ve{M}\,(T)$ & 磁化强度 \\
$m_{\rm e}^{*}\,(m_{\rm h}^{*})$ & 电子（空穴）有效质量 \\
$n$ & 粒子数 \\
$\hat{n}$ & 粒子数算符 \\
$N$ & 元胞（格点）数 \\
$\ve{p}$ & 动量 \\
$\ve{P}$ & 极化矢量 \\
$P_{k\sigma}\,(P_{k\sigma}^{*})$ & 广义动量 \\
$Q_{k\sigma}\,(Q_{k\sigma}^{*})$ & 广义坐标 \\
$\ve{q}$ & 波矢量 \\
$\ve{R}_l=\sum_{i=1}^{s}l_i\ve{a}_i\equiv\ve{l}$ & 正格矢 \\
$\hat{\ve{S}}_l=(\hat{S}_l^{\,x},\hat{S}_l^{\,y},\hat{S}_l^{\,z})$ & 自旋算符 \\
$\hat{S}_l^{+}\,(\hat{S}_l^{-})$ & 自旋上升（下降）算符 \\
$\hat{\ve{S}}\,(\ve{q})$ & 自旋密度算符 \\
$S_{\rm es}$ & 超导电子熵 \\
$t$ & 时间 \\
$T$ & 温度 \\
$T_{\rm c}$ & 临界温度 \\
$T_{\rm F}$ & 费密温度 \\
\end{tabular}
}
\par\bigbreak
% ---------- 原书第 iii 页 ----------
{\small
\begin{tabular}{@{}p{0.46\textwidth}@{}p{0.50\textwidth}@{}}
$U$ & 电子库仑排斥势能 \\
$U(T)$ & 内能 \\
$\ve{u}_l,\ \ve{u}_l(s)$ & 晶格位移 \\
$\ve{v}$ & 速度 \\
$V$ & 体积，势能 \\
$W,\ w$ & 跃迁概率 \\
$Z$ & 配分函数，配位数，离子价数 \\
$\alpha$ & 极化子的量纲为 1 的耦合常数 \\
$\beta$ & $\beta=(k_{\rm B}T)^{-1}$ \\
$\Delta$ & 能带宽度 \\
$\Delta(T)$ & 超导能隙 \\
$\Delta(\ve{r})$ & 超导能隙函数，晶体体积的相对变化 \\
$\varepsilon(\infty)\,[\varepsilon(0)]$ & 高（低）频介电常数 \\
$\varepsilon(\ve{q},\omega)$ & 介电函数 \\
$\varepsilon_{\rm k}$ & 从费密能算起的自由电子能量 \\
$\zeta(3)$ & 黎曼函数 \\
$\theta$ & 角度 \\
$\theta_{\rm D}$ & 德拜温度 \\
$\lambda$ & 波长 \\
$\lambda(T)$ & 穿透深度 \\
$\mu$ & 化学势，迁移率 \\
$\xi_0$ & BCS 相干长度 \\
$\xi_{\rm k}$ & 超导元激发能量 \\
$\rho$ & 密度，电阻率 \\
$\rho(\omega)$ & 态密度 \\
$\sigma$ & 电导率，自旋取向标号 \\
$\hat{\sigma}^{z},\ \hat{\sigma}^{+},\ \hat{\sigma}^{-}$ & 泡利矩阵 \\
$\tau$ & 寿命，弛豫时间 \\
$\Phi$ & 晶格振动位能 \\
$\Phi_0=\dfrac{2\pi\hbar c}{2e}$ & 超导磁通量子 \\
$\Phi_{\alpha\beta}(\ve{l}-\ve{l}'),\ \Phi_{\alpha\beta}\!\left(\begin{smallmatrix}\ve{l}-\ve{l}'\\ s,s'\end{smallmatrix}\right)$ & 晶格弹性力系数的张量 \\
$\phi$ & 方位角，电子波函数 \\
$\phi(\ve{r})$ & 位能密度 \\
$\Psi\,(\Psi^{+})$ & 二次量子化态向量 \\
\end{tabular}
}
\par\bigbreak
% ---------- 原书第 iv 页 ----------
{\small
\begin{tabular}{@{}p{0.46\textwidth}@{}p{0.50\textwidth}@{}}
$\psi$ & 电子波函数 \\
$\chi$ & 自旋波函数 \\
$\chi^{(k)}$ & 特征标 \\
$\chi_{\rm P}$ & 泡利磁化率 \\
$\chi^{-+}(\ve{q},\omega)$ & 横向动态磁化率 \\
$\omega$ & 频率 \\
$\omega_{\rm D}$ & 德拜频率 \\
$\omega_{\rm P}$ & 电子的等离子振荡频率 \\
$\Omega_{\rm P}$ & 离子的等离子振荡频率 \\
$\Omega$ & 正点阵元胞体积 \\
$\Omega^{*}$ & 倒点阵元胞体积 \\
${\rm d}\Omega$ & 立体角 \\
\end{tabular}
}
